\section{The complete local squareclass test}\label{sec:strata}
Write a hypothetical primitive solution as
\[
A=au^2,\qquad B=bv^2,\qquad C=cw^2,
\]
where $a,b,c$ are positive pairwise coprime squarefree integers. The equation becomes
\begin{equation}\label{eq:strata}
a^3u^6+b^5v^{10}=c^7w^{14}.
\end{equation}
At a prime $p$, a local solution is called primitive when at most one base is divisible by $p$. Roots and bases are required to be nonzero in $\Z_p$. The roots may depend on $p$. This last freedom is exactly why local existence can be weaker than an integer solution.

\begin{theorem}[Local coefficient criterion]\label{thm:localstrata}
For fixed $a,b,c$ as above, primitive local solvability of \eqref{eq:strata} is determined by the following conditions.
\begin{enumerate}
\item At two, at least one row in the table must be admissible. An even coefficient forces its own base to be the even base.
\item At an odd prime dividing $a,b,c$, respectively, the condition is
\[
\left(\frac{bc}{p}\right)=1,\qquad
\left(\frac{ac}{p}\right)=1,\qquad
\left(\frac{-ab}{p}\right)=1.
\]
\item At every odd prime not dividing $abc$, a primitive local solution exists.
\end{enumerate}
\end{theorem}
\begin{center}
\begin{tabular}{lll}\toprule
Even base&Required odd coefficients&Relation modulo eight\\\midrule
$A$&$b,c$&$b\equiv c$\\
$B$&$a,c$&$a\equiv c$\\
$C$&$a,b$&$a+b\equiv0$\\\bottomrule
\end{tabular}
\end{center}
The proof is worth retaining in full, because it establishes the absence of an untested prime power tail. The criterion is not an extrapolation from residue sweeps.

\subsection{Odd prime lifting}
For the polynomial on the left of \eqref{eq:strata} minus its right side, the partial derivatives are
\[
6a^3u^5,\qquad10b^5v^9,\qquad-14c^7w^{13}.
\]
In a primitive residue solution, at least two bases are units. At an odd prime, at most one of $6,10,14$ is divisible by $p$. Thus at least one unit variable has unit derivative. Lift the other variables to nonzero integers, using $p$ for a zero residue if necessary. Successive correction in the selected variable by multiples of $p^n$ kills one new digit at each stage. The derivative remains a unit, so the corrections converge to a nonzero root with the original primitive pattern. Every primitive residue solution therefore lifts to $\Z_p$, including at $3,5,7$.

If $p\mid a$, primitivity forces the other two bases to be units. Reduction gives $b^5v^{10}=c^7w^{14}$. The quadratic character of this equality proves the necessary symbol condition. Conversely, $v^{10}w^{-14}$ ranges over the squares of $\F_p^\times$, because $\gcd(10,14,p-1)=2$. The symbol condition lets us choose the two unit roots, then lift. The other two cases use $\gcd(6,14,p-1)=2$ and $\gcd(6,10,p-1)=2$. Moving a term across the equality gives the minus sign when $p\mid c$.

\subsection{Odd primes away from the coefficients}
Let $\alpha,\beta,\gamma$ be the quadratic characters of $a,b,c$, and let $e=(-1/p)$. A primitive residue solution with one zero base exists if
\[
\beta=\gamma,\qquad\alpha=\gamma,\qquad\alpha=e\beta,
\]
respectively. The same cyclic group argument proves this, and the lifting construction replaces the zero root by a nonzero multiple of $p$.

If $p\equiv1\pmod4$, two of the three signs agree, so one of these conditions holds. If $p\equiv3\pmod4$ and all three fail, then $\alpha=\beta=-\gamma$. The two remaining sign patterns can be handled by the universal identity of Proposition~\ref{prop:universal}.

There are $(p+1)/4$ elements $t$ for which both $t$ and $1-t$ are nonsquares. Here is a direct count. Extend the quadratic character $\chi_p$ by zero. Counting $(x-y)(x+y)=d$ for $d\ne0$ gives
\[
\sum_x\chi_p(x^2-d)=-1.
\]
Completing the square then gives $\sum_t\chi_p(t(1-t))=-\chi_p(-1)=1$. After omitting $0,1$, both single character sums are $-1$. Expanding the indicator $(1-\chi_p(t))(1-\chi_p(1-t))/4$ over the remaining points gives $(p+1)/4>0$.

For such a $t$, the three bases in Proposition~\ref{prop:universal} have characters $(1,1,-1)$. Multiplication by the weighted powers of a nonsquare $\lambda$ changes all three signs and gives $(-1,-1,1)$. In either needed sign pattern, the ratios of these bases to $a,b,c$ are squares. Choosing their square roots gives a solution of \eqref{eq:strata} with all bases units. The lifting argument completes the proof at these primes.

\subsection{The dyadic criterion}
Exactly one base is even in a primitive two adic solution. An even powered base is zero modulo eight, while an odd squared root is one modulo eight and an odd coefficient has the same third, fifth, and seventh power as itself modulo eight. This proves the necessity of the table.

For sufficiency, fix the selected even root to be one when its coefficient is even, and two when the coefficient is odd. Fix one odd root to be one. After moving the terms and dividing by the remaining coefficient power, the required remaining root power is a unit congruent to one modulo eight. Such a unit is a square in $\Z_2$, and every odd power map on $\Z_2^\times$ is bijective. It consequently has the required sixth, tenth, or fourteenth root.

The square assertion has a direct lifting proof. For $n\ge3$, changing an odd trial square root by $2^{n-1}$ changes its square by $2^n$ modulo $2^{n+1}$, allowing the next digit to be selected. Odd power maps have odd derivative and lift uniquely. Thus the dyadic criterion also covers every higher precision.

\subsection{Two families that survive every local refinement}
The coefficients $(a,b,c)=(3,5,2)$ pass the even $C$ row because $3+5\equiv0\pmod8$. At three, $bc=10$ is a square residue. At five, $ac=6$ is a square residue. Hence
\[
27u^6+3125v^{10}=128w^{14}
\]
has a primitive solution at every prime. Likewise $(2,3,11)$ passes the even $A$ row since $3\equiv11\pmod8$. At three, $ac=22\equiv1$, and at eleven, $-ab=-6\equiv5=4^2$. Thus
\[
8u^6+243v^{10}=11^7w^{14}
\]
also has a primitive solution at every prime. Both have positive real solutions. Neither is asserted to have an integer solution.

For either coefficient triple and any modulus $M$, there are positive integer roots giving pairwise coprime bases and satisfying the equation modulo $M$. Enlarge $M$ to contain the coefficient primes, choose primitive local residues, and combine them by the Chinese remainder theorem. Choose $u$ in its positive progression. Choose $v$ in its progression with $v\equiv1$ at each prime dividing $u$ outside the enlarged modulus. Choose $w$ with the analogous conditions at primes dividing $uv$. The auxiliary moduli are coprime to the original one. Inside the modulus, local primitivity prevents shared base primes. Outside it, the avoidance conditions do so. This proves the finite congruence assertion, while exact equality remains unproved.

Joint labels still recover real exclusions. Among the 64 all odd coefficient triples modulo eight, the three row criterion allows 36. Every individual coordinate marginal is full. Replacing the joint relation by independent marginals restores 28 excluded triples. Packet refinement therefore has two distinct lessons: preserve correlations that already contain arithmetic information, and recognize when the projected local problem has become completely soluble.
